\documentclass[12pt]{article}

% --- Packages ---
\usepackage[utf8]{inputenc}
\usepackage{geometry}
\geometry{a4paper, margin=1in}
\usepackage{xcolor}
\usepackage{listings}

% --- Code Highlighting Setup ---
\definecolor{codegreen}{rgb}{0,0.6,0}
\definecolor{codegray}{rgb}{0.5,0.5,0.5}
\definecolor{codepurple}{rgb}{0.58,0,0.82}
\definecolor{backcolour}{rgb}{0.95,0.95,0.92}

\lstset{
    backgroundcolor=\color{backcolour},   
    commentstyle=\color{codegreen},
    keywordstyle=\color{magenta},
    numberstyle=\tiny\color{codegray},
    stringstyle=\color{codepurple},
    basicstyle=\ttfamily\small,
    breaklines=true,
    tabsize=2
}

% --- Document Metadata ---
\title{Operational Execution Logic \\ \large Machine-Level Implementation}
\date{}

\begin{document}

\maketitle

\section{Memory Allocation and Pointers}
Instantiate acts as a memory allocation command that requests a specific block of heap space from the system. When this operation is executed, the runtime calculates the byte-size of the class template, carves out that space in RAM, and returns a memory address. This address serves as a pointer, allowing the processor to locate and interact with the data stored at that specific physical coordinate in the hardware.

\begin{lstlisting}[language=Java]
// System requests heap space and returns memory address
Guitar myGuitar = new Guitar(); 
\end{lstlisting}

\section{Bitwise Movement and Gating}
Assign is the mechanical movement of bits from one register to a designated memory offset. From a machine perspective, it overwrites the existing electrical states at a specific address with new data. When gated by methods, the CPU must first execute a validation check, ensuring the incoming bit pattern is compatible before allowing the write operation to modify the object's internal state in memory.

\begin{lstlisting}[language=Java]
// Overwriting electrical states at a designated memory offset
myGuitar.setVolume(11); 
\end{lstlisting}

\section{Control Flow and Instruction Pointers}
Compare and Loop function as hardware-level control flow instructions that manipulate the instruction pointer. A comparison causes the CPU to evaluate the difference between two values and set status flags in its register. If the resulting flag matches a specific condition, the system executes a jump instruction that resets the instruction pointer to a previous memory address, forcing the processor to repeat a sequence of operations until the flag state changes.

\begin{lstlisting}[language=Java]
// Evaluating values to set CPU flags and jump instruction pointers
for (Guitar g : studioRack) {
    if (g.getVolume() > 10) {
        g.tune();
    }
}
\end{lstlisting}

\end{document}
